% !TeX root = ic-screen.tex

% ic.tex
%
% driver file ic.tex to produce text

\preto\OLEndChapterHook{\IfFileExists{include/summary-\thechapter}
        {\section*{摘要}\addcontentsline{toc}{section}{摘要}
        \let\emph\textbf\input{include/summary-\thechapter}\let\emph\textit}{}}

\problemsperchapter
\allowdisplaybreaks

\frontmatter

\OLPfrontmatter

% !TeX root = ../ic-screen.tex

\chapter{前言}

Gödel（哥德尔）不完全性定理是数理逻辑——即使不是整个数学——中最著名的成果之一。
其第一定理粗略地说是指，任何能够执行最少量算术的可公理化数学理论，要么是不一致的（平凡的，能证明一切），要么是不完全的。
换言之，如果一个数学理论能够以紧凑的方式写下来（即可公理化），其强度足以陈述并证明关于自然数的一些基本事实，并且不包含使其失效的矛盾，那么在该理论的语言中，总存在一些它无法确定的陈述。
也就是说，存在语句~$A$，使得该理论既不能证明$A$，也不能证明~$\lnot A$。（这是第一不完全性定理。）

这一结果在历史上令人惊讶，因为它可能意味着我们永远无法写出一个能够“捕获”所有数学真理的公理理论——至少如果我们要求从理论得出的推论必须能通过有限的推导或证明来验证，而这些推导或证明的正确性又可被机械地检验。
该结果的一个推论是，数学真理是不可判定的：不存在一种机械的方法，能够判定一个数学理论中的语句是否可以从公理推出。
另一个推论是，足够强且一致的数学理论无法\emph{证明}其自身的一致性——至少当以最直接的方式在理论内部形式化一致性陈述时是如此。（这是第二不完全性定理。）

假设读者具备形式逻辑的基本背景，实际上陈述并证明这些结果中的前两个并不困难。
事实上，我们在下面的\cref{inc:int::chap}中就会这样做。
但我们在此证明的版本对第一不完全性定理作了比实际所需更强的假设。
我们也没有证明这些假设对我们所讨论的理论成立，并且证明是非构造性的：它们证明了理论是不完全的，但并未给出理论无法确定的语句的实例。
要详细地做到这一点需要更多背景知识。

在\cref{cmp:rec::chap,inc:art::chap,inc:req::chap}中，我们提供这些背景知识：我们讨论一种可计算性模型（原始递归函数），说明通过给符号指派数字，我们可以用原始递归函数和关系来“算术化”算术理论的语法和证明论，最后证明非常简单的算术理论~$\Th{Q}$能够“表示”所有的原始递归函数和关系。

有了这些背景，就有可能以与任何关于这些结果的详尽数学阐述相同的详细程度来陈述和证明不完全性定理。
在\cref{inc:inp::chap}中，我们证明了Gödel第一不完全性定理的原始版本、Rosser（罗瑟）的改进版本、第二不完全性定理，以及Löb定理和Tarski（塔斯基）关于真的不可定义性的定理。

不完全性现象与可计算性的概念紧密相连。
不完全性定理适用于那些能够从可计算的公理集可计算地生成的理论。
语法的算术化的机制需要一个计算模型（递归函数）来阐明细节。
而该结果本身在结构上与可计算性理论中的一个著名结果相关，即由Church（邱奇）和Turing（图灵）证明的停机问题不可判定这一定理。
因此，很自然地可以从可计算性理论的角度，给出另一种关于不完全性定理的描述和证明。
我们在引入部分可计算函数以及可判定集与可计算枚举集的基本理论之后，在\cref{comp-inc:chap}中这样做。

算术理论是不完全的，这意味着它们不仅拥有一些不像自然数“标准模型”的模型，而且还会使得一些在标准模型中为真的语句在这些模型中为假。
这类模型的结构本身就是一个有趣的研究领域。
我们在\cref{mod:chap}中对此作简要的探讨。

不完全性定理首要关注的是那些用一阶语言表述并可公理化的理论，并且假定它们具有通常的一阶后承关系与可证性关系。
根据Gödel完备性定理，我们知道我们为一阶逻辑及建立于其上的理论所配备的证明系统达到了我们所希望的强度：它们能推出公理的所有后承。
在某种意义上，不完全性定理指出，恰恰是一阶逻辑的这种特性，阻碍了我们写下能“解决每个问题”的数学理论。
一个自然的出路是采用一种具有更强后承关系的更强逻辑：二阶逻辑。
我们将在\cref{sol:chap}中讨论它最重要的性质：它强到足以刻画算术真，并且其表达力远超一阶逻辑。
二阶算术是完备的。但是，许多对一阶逻辑成立的结果（如紧致性定理和Löwenheim--Skolem定理）对二阶逻辑并不成立。

本书明确引入递归函数作为可计算性的一个模型。
$\Th{Q}$中的可表示性也可以被视为可计算性的一个模型，并且在此过程中，我们（几乎）证明了它与递归函数等价。
配套书籍《\href{https://slc.openlogicproject.org/}{\emph{Sets, Logic, Computation}}》讨论了Turing机这一计算模型。
在\cref{lambda:chap}中，我们还将引入另一种可计算性模型：（无类型的）λ演算。

\section*{给教师的说明}

这是一本关于Gödel不完全性定理和递归函数论的教科书。
我在卡尔加里大学教授哲学479（逻辑学III）时将其作为主要教材。
它基于\href{https://openlogicproject.org}{开放逻辑项目}的材料。

顾名思义，这门课是一个系列中的第三门，因此学生（也就是本书的读者）应该已经熟悉一阶逻辑。
（逻辑学I使用教材《\href{https://forallx.openlogicproject.org}{\emph{forall x: Calgary}}》，
逻辑学II使用另一本基于OLP的教材《\href{https://slc.openlogicproject.org}{\emph{Sets, Logic, Computation}}》。）
不过，逻辑学II所要求的知识已收录在\cref{fol:chap,nd:chap}中，
因此，基于本书授课时，并不绝对需要比此更多的预备知识。

逻辑学III是一门十三周的课程，每周三课时。
这通常足以覆盖\cref{inc:int::chap,cmp:rec::chap,inc:art::chap,inc:req::chap,inc:inp::chap}以及\cref{comp-inc:chap,mod:chap,sol:chap,lambda:chap}中的一两个专题，
具体取决于学生的兴趣。如果学生对一阶逻辑的基础知识，特别是自然演绎还不熟悉，
你可能需要在这些内容上花更多时间。
请注意，在正文中讨论算术理论（如$\Th{Q}$和$\Th{PA}$）中的可证性时，
可证性断言的证明并不是使用某个特定的推导系统给出的。
相反，我们只是直观地论证了某些断言可以依据一阶逻辑从公理推出。
然而，\cref{deriv:chap}包含了大量从~$\Th{Q}$的公理出发进行自然演绎推导的实际例子。
\Cref{inc:art::chap}则为自然演绎进行语法的算术化。
这或许是本书的一个独特之处；大多数其他教科书只对公理化推导进行算术化。
后者的编码要容易得多，但用它来给出证明则困难得多。

\section*{致谢}

\cref{inc:int::chap,cmp:rec::chap,inc:art::chap,inc:req::chap,inc:inp::chap,comp-inc:chap,lambda:chap}中所使用的OLP材料最初基于Jeremy Avigad贡献给OLP的
“可计算性与不完全性”讲义。我对这些材料进行了大量修订和扩充。
例如，原来的讲义基于公理化推导系统来构建算术理论。
而这里我们使用Gentzen的标准自然演绎系统（在\cref{nd:chap}中描述），
这需要用原始递归的方式处理树（在\cref{cmp:rec:tre:sec}中），
并采用更复杂的方法对推导进行算术化（在\cref{inc:art:pnd:sec}中）。
\cref{lambda:chap}中的材料也在钱泽森作为Mitacs暑期实习生访问卡尔加里期间由他进行了扩充。

OLP中关于模型论和算术模型的材料（在\cref{mod:chap}中）
最初取自Aldo Antonelli贡献给OLP的讲义“经典命题与谓词逻辑的完备性”。
他在2015年不幸早逝。

\cref{bios:chap}中的逻辑学家传记以及\cref{nd:chap}中的大部分材料最初由Sam Burns完成。
Dana Hägg最初参与了\cref{fol:chap}材料的编写工作。

\mainmatter

\olimport*[incompleteness/introduction]{introduction}

\olimport*[computability/recursive-functions]{recursive-functions}

\olimport*[incompleteness/arithmetization-syntax]{arithmetization-syntax}

\olimport*[incompleteness/representability-in-q]{representability-in-q}

\olimport*[incompleteness/incompleteness-provability]{incompleteness-provability}

\chapter{可计算性与不完全性}\label{comp-inc:chap}

\olimport*[computability/computability-theory]{introduction}
\olimport*[computability/computability-theory]{coding-computations}
\olimport*[computability/computability-theory]{normal-form}
\olimport*[computability/computability-theory]{s-m-n}
\olimport*[computability/computability-theory]{universal-part-function}
\olimport*[computability/computability-theory]{no-universal-function}
\olimport*[computability/computability-theory]{halting-problem}
%\olimport*[computability/computability-theory]{russells-paradox}
\olimport*[computability/computability-theory]{computable-sets}
\olimport*[computability/computability-theory]{ce-sets}
\olimport*[computability/computability-theory]{equiv-ce-defs}
\olimport*[computability/computability-theory]{non-comp-set}
\olimport*[computability/computability-theory]{complement-ce}
\olimport[include]{ce-theories}
\olimport[include]{g1-halting-problem}

\OLEndChapterHook

\chapter{算术模型}\label{mod:chap}

\olimport*[model-theory/models-of-arithmetic]{introduction}

\olimport*[model-theory/basics]{reducts-and-expansions}

\olimport*[model-theory/basics]{isomorphism}

\olimport*[model-theory/basics]{theory-of-m}

\olimport*[model-theory/models-of-arithmetic]{standard-models}

\olimport*[model-theory/models-of-arithmetic]{non-standard-models}

\olimport*[model-theory/models-of-arithmetic]{models-of-q}

\olimport*[model-theory/models-of-arithmetic]{models-of-pa}

\olimport*[model-theory/models-of-arithmetic]{computable-models}

\OLEndChapterHook

\chapter{二阶逻辑}\label{sol:chap}

\olimport*[second-order-logic/syntax-and-semantics]{introduction}

\let\oldolsection\olsection
\let\olsection\nosection
\olimport*[second-order-logic/sol-and-set-theory]{introduction}
\let\olsection\oldolsection

\olimport*[second-order-logic/syntax-and-semantics]{terms-formulas}

\olimport*[second-order-logic/syntax-and-semantics]{satisfaction}

\olimport*[second-order-logic/syntax-and-semantics]{semantic-notions}

\olimport*[second-order-logic/syntax-and-semantics]{expressive-power}

\olimport*[second-order-logic/syntax-and-semantics]{inf-count}

\olimport*[second-order-logic/metatheory]{second-order-arithmetic}

\olimport*[second-order-logic/metatheory]{undecidability-and-axiomatizability}

\olimport*[second-order-logic/metatheory]{compactness}

\olimport*[second-order-logic/metatheory]{loewenheim-skolem}

\olimport*[second-order-logic/sol-and-set-theory]{comparing-sets}

\olimport*[second-order-logic/sol-and-set-theory]{cardinalities}

\olimport*[second-order-logic/sol-and-set-theory]{power-of-continuum}

\OLEndChapterHook

\chapter{λ演算}\label{lambda:chap}

\olimport*[lambda-calculus/introduction]{overview}

\olimport*[lambda-calculus/introduction]{syntax}

\olimport*[lambda-calculus/introduction]{reduction}

\olimport*[lambda-calculus/introduction]{church-rosser}

\olimport*[lambda-calculus/introduction]{currying}

\olsection{λ可定义性}

\let\oldolsection\olsection
\def\olsection#1{}
\olimport*[lambda-calculus/lambda-definability]{introduction}
\let\olsection\oldolsection

\olimport*[lambda-calculus/lambda-definability]{arithmetical-functions}
\olimport*[lambda-calculus/lambda-definability]{pairs}
\olimport*[lambda-calculus/lambda-definability]{truth-values}
%\olimport{lists}
\olimport*[lambda-calculus/lambda-definability]{primitive-recursive-functions}
\olimport*[lambda-calculus/lambda-definability]{fixpoints}
\olimport*[lambda-calculus/lambda-definability]{minimization}
\olimport*[lambda-calculus/lambda-definability]{partial-recursive-functions}
\olimport*[lambda-calculus/lambda-definability]{lambda-definable-recursive}

\OLEndChapterHook

\appendix

\chapter{算术理论中的推导}\label{deriv:chap}

在我们证明所有一般递归函数都在~$\Th{Q}$ 中可表示，以及证明不完全性定理时，我们曾断言各种结论在 $\Th{Q}$ 和~$\Th{PA}$ 中是可证的。然而，这些断言的证明只是非形式地给出了论证，并未展示自然演绎中的实际推导。我们将在本章提供其中一些推导。

例如，在 \olref[inc][req][bre]{lem:q-proves-add} 中，我们通过对 $m$ 归纳，证明了对于所有 $n$ 和 $m \in \Nat$，$\Th{Q} \Proves (\num{n} +
\num{m}) = \num{n+m}$。

\begin{proof}[{\olref[inc][req][bre]{lem:q-proves-add}} 的证明]
基始：$m = 0$。此时需要证明的是，对所有 $n$，$\Th{Q} \Proves \num{n} + \num{0} = \num{n+0}$。由于 $\num{0}$ 就是 $\Obj 0$ 而 $\num{n+0}$ 就是 $\num{n}$，这就相当于要证明$\Th{Q} \Proves (\num{n} + \Obj 0) = \num{n}$。推导
\begin{prooftree}
  \AxiomC{$\lforall[x][(x + \Obj 0) = x]$}
  \RightLabel{\Elim\lforall}
  \UnaryInfC{$(\num{n} + \Obj 0) = \num{n}$}
\end{prooftree}
是$(\num{n} + \Obj 0) = \num{n}$ 的一个自然演绎推导，它有一个未解除假设，而这个未解除假设正是~$\Th{Q}$ 的一条公理。

归纳步骤：设对任意 $n$，已有$\Th{Q} \Proves (\num{n} +
\num{m}) = \num{n+m}$（例如，通过推导 $\delta_{n,m}$）。还需证明$\Th{Q} \Proves (\num{n} + \num{m+1}) =
\num{n+m+1}$。注意到$\num{m+1} \ident \num{m}'$，且$\num{n+m+1} \ident \num{n+m}'$。因此，我们寻求从~$\Th{Q}$ 的公理出发，推导出$(\num{n} + \num{m}') = \num{n+m}'$。我们的推导可以使用归纳假设所保证存在的推导 $\delta_{n,m}$。
\begin{prooftree}
  \AxiomC{}
  \RightLabel{$\delta_{n,m}$}
  \DeduceC{$(\num{n} + \num{m}) = \num{n+m}$}
  \AxiomC{$\lforall[x][\lforall[y][(x+y') = (x+y)']]$}
  \RightLabel{\Elim\lforall}
  \UnaryInfC{$\lforall[y][(\num{n}+y') = (\num{n}+y)']$}
  \RightLabel{\Elim\lforall}
  \UnaryInfC{$(\num{n}+\num{m}') = (\num{n}+\num{m})'$}
    \RightLabel{\Elim=}
    \BinaryInfC{$(\num{n}+\num{m}') = \num{n+m}'$}
\end{prooftree}
在最后的 $\Elim=$ 推理中，我们将右前提右端$(\num{n} + \num{m})'$ 中的子项 $\num{n} + \num{m}$ 替换为项 $\num{n+m}$。
\end{proof}

在 \olref[inc][req][min]{lem:less-zero} 中，我们证明了$\Th{Q} \Proves
\lforall[x][\lnot x < \Obj 0]$。一个实际的推导会是什么样子呢？

\begin{proof}[{\olref[inc][req][min]{lem:less-zero}} 的证明]
为了证明这种全称结论，我们使用 $\Intro\lforall$，它需要先推导出 $\lnot a < \Obj 0$。考察公理 $!Q_8$，这意味着要证明 $\lnot \exists z (z' + a) = \Obj 0$。具体来说，若能证明后者，则由 $!Q_8$ 即可证明前者（回想一下，$A \liff B$ 是 $(A \lif B) \land (B \lif
A)$ 的缩写）。
\begin{prooftree}\footnotesize
  \AxiomC{$\lnot\lexists[z][(z' + a) = \Obj 0]$}
  \AxiomC{$\lforall[x][\lforall[y][(x < y \liff \lexists[z][(z' + x) = y])]]$}
  \RightLabel{\Elim\lforall}
  \UnaryInfC{$\lforall[y][(a < y \liff \lexists[z][(z' + a) = y])]$}
  \RightLabel{\Elim\lforall}
  \UnaryInfC{$a < \Obj 0 \liff \lexists[z][(z' + a) = \Obj 0]$}
  \RightLabel{\Elim\land}
  \UnaryInfC{$a < \Obj 0 \lif \lexists[z][(z' + a) = \Obj 0]$}
  \AxiomC{$\Discharge{a < \Obj 0}{1}$}
  \RightLabel{\Elim\lif}
  \BinaryInfC{$\lexists[z][(z' + a) = \Obj 0]$}
  \RightLabel{\Elim\lnot}
  \insertBetweenHyps{\hskip -3em}
  \BinaryInfC{$\lfalse$}
  \DischargeRule{\Intro\lnot}{1}
  \UnaryInfC{$\lnot a<\Obj 0$}
\end{prooftree}
这是从 $\lnot\lexists[z][(z' +
a) = \Obj 0]$（以及公理~$!Q_8$）出发的 $\lnot a<\Obj 0$ 的推导；我们称它为~$\delta_1$。

现在，我们如何从~$\Th{Q}$ 的公理证明 $\lnot\lexists[z][(z' + a) = \Obj 0]$？要证明像这样的否定式，我们需要如下形式的推导：
\begin{prooftree}
  \AxiomC{$\Discharge{\lexists[z][(z' + a) = \Obj 0]}{2}$}
  \DeduceC{$\lfalse$}
  \DischargeRule{\Intro\lnot}{2}
  \UnaryInfC{$\lnot\lexists[z][(z' + a) = \Obj 0]$}
  \end{prooftree}
为了从一个存在性断言得到矛盾，我们为存在量化的变元~$z$ 引入一个常项~$b$，并使用 \Elim\lexists：
\begin{prooftree}
  \AxiomC{$\Discharge{\lexists[z][(z' + a) = \Obj 0]}{2}$}
  \AxiomC{$\Discharge{(b'+a) = \Obj 0}{3}$}
  \RightLabel{$\delta_2$}
    \DeduceC{$\lfalse$}
    \DischargeRule{\Elim\exists}{3}
    \BinaryInfC{$\lfalse$}
  \DischargeRule{\Intro\lnot}{2}
  \UnaryInfC{$\lnot\lexists[z][(z' + a) = \Obj 0]$}
  \end{prooftree}
现在的任务是填充 $\delta_2$，即从 $(b'+a) = \Obj 0$ 和~$\Th{Q}$ 的公理证明 $\lfalse$。$Q_2$ 说 $\Obj 0$ 不可能是某个数的后继，所以一种方法是证明 $(b' + a)$ 等于某个数的后继。由于这个表达式本身是一个和式，加法的公理必须发挥作用。若 $\eq[a][\Obj 0]$，则由 $Q_5$ 可得 $\eq[(b'
+ a)][b']$，也就是说 $b' + a$ 是某个数（即~$b$）的后继。另一方面，如果对某个 $c$ 有 $\eq[a][c']$，那么由 \Elim\eq 可得 $\eq[(b'+a)][(b'+c')]$，再由~$Q_6$ 可得 $\eq[(b'+c')][(b'+c)']$。因此 $b'+a$ 同样是某个数的后继——这次是 $b'+c$。因此，策略是将任务分为这两种情况。我们还必须验证 $\Th{Q}$ 能证明其中一种情况成立，即 $\Th{Q} \Proves a = 0 \lor \lexists[y][(a = y')]$，但这直接由 $Q_3$ 通过 \Elim\lforall 得出。以下是两种情况：

情况 1：从 $\eq[a][\Obj 0]$ 和 $\eq[(b'+a)][\Obj
  0]$（以及公理 $Q_2$、$Q_5$）证明 $\lfalse$：
\begin{prooftree}\footnotesize
  \AxiomC{$\lforall[x][\lnot \Obj 0 = x']$}
  \RightLabel{\Elim\lforall}
  \UnaryInfC{$\lnot \Obj 0 = b'$}
  \AxiomC{$\lforall[x][(x+\Obj 0) = x]$}
  \RightLabel{\Elim\lforall}
  \UnaryInfC{$(b' + \Obj 0) = b'$}
  \AxiomC{$a = \Obj 0$}
  \AxiomC{$(b'+a) = \Obj 0$}
  \RightLabel{\Elim=}
  \BinaryInfC{$(b' + \Obj 0) = \Obj 0$}
  \doubleLine
  \UnaryInfC{$\Obj 0 = (b' + \Obj 0)$}
      \insertBetweenHyps{\hskip -.5em}
  \RightLabel{\Elim=}
  \BinaryInfC{$\Obj 0 = b'$}
  \RightLabel{\Elim\lnot}
  \BinaryInfC{$\lfalse$}
\end{prooftree}
称这个推导为~$\delta_3$。（我们通过一条双推理线，简写了从 $(b' + \Obj 0) = \Obj 0$ 到 $\Obj 0 = (b' + \Obj 0)$ 的推导。）

情况 2：从 $\lexists[y][a = y']$ 和 $\eq[(b'+a)][\Obj 0]$（以及公理 $Q_2$、$Q_6$）证明 $\lfalse$。我们首先展示如何从 $\eq[a][c']$ 和 $\eq[(b'+a)][\Obj 0]$ 推导出 $\lfalse$。
\begin{prooftree}\footnotesize
  \AxiomC{$\lforall[x][\lnot \Obj 0 = x']$}
  \RightLabel{\Elim\lforall}
  \UnaryInfC{$\lnot \Obj 0 = (b'+c)'$}
  \AxiomC{$a = c'$}
  \AxiomC{$(b'+a) = \Obj 0$}
  \RightLabel{\Elim=}
          \insertBetweenHyps{\hskip -.3em}
  \BinaryInfC{$\Obj (b'+c') = \Obj 0$}
  \AxiomC{$\lforall[x][\lforall[y][(x+y') = (x+y)']]$}
  \RightLabel{\Elim\lforall}
  \UnaryInfC{$\lforall[y][(b'+y') = (b'+y)']$}
  \RightLabel{\Elim\lforall}
  \UnaryInfC{$(b' + c') = (b'+c)'$}
  \RightLabel{\Elim=}
        \insertBetweenHyps{\hskip .5em}
  \BinaryInfC{$\Obj 0 = (b' + c)'$}
  \RightLabel{\Elim\lnot}
          \insertBetweenHyps{\hskip -.5em}
  \BinaryInfC{$\lfalse$}
\end{prooftree}
称此推导为 $\delta_4$。我们通过应用 $\Elim\lexists$ 并解除假设 $\eq[a][c']$，得到所需的推导 $\delta_5$：
\begin{prooftree}
  \AxiomC{$\lexists[y][a=y']$}
  \AxiomC{$\Discharge{a = c'}{6} \quad \eq[(b'+a)][\Obj 0]$}
  \RightLabel{$\delta_4$}
  \DeduceC{$\lfalse$}
  \DischargeRule{\Elim\exists}{6}
  \BinaryInfC{$\lfalse$}
\end{prooftree}

  
将所有部分组合起来，完整的证明如下所示：
\begin{prooftree}\footnotesize
  \AxiomC{$\Discharge{\lexists[z][(z' + a) = \Obj 0]}{2}$}
  \AxiomC{$\begin{gathered}
  \lforall[x][(x = 0 \lor {}]\\
  \lexists[y][(a = y')])
  \end{gathered}$}
  \RightLabel{\Elim\lforall}
  \UnaryInfC{$\begin{gathered}a = 0 \lor {}\\
  \lexists[y][(a = y')]
  \end{gathered}$}
  \AxiomC{$\begin{gathered}[b]
      \Discharge{a = \Obj 0}{7} \\
      \Discharge{(b'+a) = \Obj 0}{3}
      \end{gathered}$}
    \RightLabel{$\delta_3$}
    \DeduceC{$\lfalse$}
    \AxiomC{$\begin{gathered}[b]
        \Discharge{\lexists[y][a=y']}{7} \\
        \Discharge{(b'+a) = \Obj 0}{3}
        \end{gathered}$\quad}
    \RightLabel{$\delta_5$}
    \DeduceC{$\lfalse$}
    \DischargeRule{\Elim\lor}{7}
    \insertBetweenHyps{\hskip -.5em}
    \TrinaryInfC{$\lfalse$}
    \RightSubproofLabel{$\delta_2$}
    \DischargeRule{\Elim\exists}{3}
    \BinaryInfC{$\lfalse$}
  \DischargeRule{\Intro\lnot}{2}
  \UnaryInfC{$\lnot\lexists[z][(z' + a) = \Obj 0]$}
  \RightLabel{$\delta_1$}
  \DeduceC{$\lnot a<\Obj 0$}
  \RightLabel{\Intro\lforall}
  \UnaryInfC{$\lforall[x][\lnot x < \Obj 0]$}
\end{prooftree}
\end{proof}

在 \olref[inc][inp][ros]{thm:rosser} 的证明中，我们将
$\ORProv(y)$ 定义为 \[\lexists[x][(\OPrf(x, y) \land \lforall[z][(z < x
    \lif \lnot \ORefut(z, y))])].\]其中
$\OPrf(x,y)$ 是在 $\Th{Q}$ 中表示理论~$\Th{T}$（一个一致且可公理化的 $\Th{Q}$ 的扩张）的证明关系的公式，而 $\ORefut(z, y)$ 是表示反驳关系的公式。
这意味着，如果 $n$ 是~$!A$ 的一个证明的Gödel数，那么 $\Th{Q} \Proves
\OPrf(\num{n}, \gn{!A})$，否则 $\Th{Q} \Proves
\lnot\OPrf(\num{n}, \gn{!A})$。
类似地，如果 $n$ 是 $\lnot !A$ 的一个证明的Gödel数，那么 $\Th{Q} \Proves \ORefut(\num{n},
\gn{!A})$，否则 $\Th{Q} \Proves \lnot\ORefut(\num{n},
\gn{!A})$。
我们利用对角线引理找到一个语句 $!R$，使得 $\Th{Q} \Proves !R \liff \lnot \ORProv(\gn{!R})$。
Rosser定理断言 $\Th{T} \Proves/ !R$ 且 $\Th{T} \Proves/ \lnot !R$。
这两个结论都是间接证明的：我们证明，若 $\Th{T} \Proves !R$，则 $\Th{T}$ 不一致，即 $\Th{T} \Proves \lfalse$；若 $\Th{T} \Proves \lnot !R$，同样如此。

\begin{proof}[Proof of {\olref[inc][inp][ros]{thm:rosser}}]
首先，我们证明关于~$<$ 的一些事实。由
\olref[inc][req][min]{lem:less-nsucc}，我们知道对每个~$n$，$\Th{Q} \Proves
\lforall[x][(x < \num {n+1} \lif (\eq[x][\Obj 0] \lor \dots \lor
  \eq[x][\num n]))]$。因此当然（若 $n>1$）也有
$\Th{Q} \Proves \lforall[x][(x < \num {n} \lif (\eq[x][\Obj 0] \lor
  \dots \lor \eq[x][\num{n-1}]))]$。我们可以利用这一点从 $a < \num{n}$ 推导出
$\eq[a][\Obj 0] \lor \dots \lor \eq[a][\num{n-1}]$：
\begin{prooftree}
  \AxiomC{$a < \num{n}$}
  \AxiomC{}
  \DeduceC{$\lforall[x][(x < \num{n} \lif (x = \num{0} \lor 
      \dots \lor x = \num{n-1}))]$}
  \RightLabel{\Elim\forall}
  \UnaryInfC{$a < \num{n} \lif (a = \num{0} \lor 
      \dots \lor a = \num{n-1})$}
  \RightLabel{\Elim\lif}
  \BinaryInfC{$a = \num{0} \lor \dots \lor a = \num{n-1}$}
\end{prooftree}
我们称这个推导为 $\lambda_1$。

现在，为证明 $\Th{T} \Proves/ !R$，我们假设 $\Th{T}
\Proves !R$（有一个推导~$\delta$），并证明这将导致 $\Th{T}$ 不一致。令 $n$ 为~$\delta$ 的Gödel数。由于 $\OPrf$ 在~$\Th{Q}$ 中表示证明关系，所以存在 $\OPrf(\num{n},
\gn{!R})$ 的一个推导~$\delta_1$。此外，由于假定 $\Th{T}$ 是一致的，没有 $k < n$ 是~$!R$ 的反驳的Gödel数，因此对每个 $k < n$，都有 $\Th{Q} \Proves \lnot \ORefut(\num{k}, \gn{!R})$；令
$\rho_k$ 为相应的推导。我们得到 $\ORProv(\gn{!R})$ 的一个推导：
\begin{prooftree}\footnotesize
  \AxiomC{}
  \RightLabel{$\delta_1$}
  \DeduceC{$\OPrf(\num{n}, \gn{!R})$}

  \AxiomC{$\Discharge{a < \num{n}}{1}$}
  \RightLabel{$\lambda_1$}
  \DeduceC{$\begin{gathered}[b]
      a= \num{0} \lor \dots {} \\
      {} \lor a = \num{n-1}
      \end{gathered}$}
  \AxiomC{$\dots$}

  \AxiomC{$\Discharge{a=\num{k}}{2}$}
  \AxiomC{}
  \RightLabel{$\rho_k$}
  \DeduceC{$\lnot \ORefut(\num{k}, \gn{!R})$}
  \RightLabel{\Elim=}
  \BinaryInfC{$\lnot \ORefut(a, \gn{!R})$}
  \AxiomC{$\dots$}
  \DischargeRule{$\Elim\lor^*$}{2}
  \doubleLine
  \insertBetweenHyps{\hskip -1pt}
  \QuaternaryInfC{$\lnot \ORefut(a, \gn{!R})$}
  \DischargeRule{\Intro\lif}{1}
  \UnaryInfC{$a < \num{n} \lif \lnot \ORefut(a, \gn{!R})$}
  \RightLabel{\Intro\lforall}
  \UnaryInfC{$\lforall[z][(z < \num{n} \lif \lnot \ORefut(z, \gn{!R}))]$}

    \insertBetweenHyps{\hskip -5pt}
  \RightLabel{\Intro\land}
  \BinaryInfC{$\OPrf(\num{n}, \gn{!R}) \land \lforall[z][(z < \num{n} \lif \lnot \ORefut(\num{z}, \gn{!R}))]$}
  \RightLabel{\Intro\lexists}
  \UnaryInfC{$\lexists[x][(\OPrf(x, \gn{!R}) \land \lforall[z][(z < x \lif \lnot \ORefut(z, \gn{!R}))])]$}
\end{prooftree}
（上面我们用 $\Elim\lor^*$ 来缩写多次应用 $\Elim\lor$。）我们已经证明了，如果 $\Th{T} \Proves !R$，则存在~$\ORProv(\gn{!R})$ 的一个推导。那么，由于 $\Th{T} \Proves R \liff
\lnot \ORProv(\gn{!R})$，也有 $\Th{T} \Proves \ORProv(\gn{!R}) \lif
\lnot R$，我们将有 $\Th{T} \Proves \lnot !R$，从而 $\Th{T}$ 不一致。

现在证明 $\Th{T} \Proves/ \lnot !R$。同样，假设如此。那么就存在 $\lnot !R$ 的一个推导 $\rho$，其Gödel数为 $m$——即~$!R$ 的一个反驳——从而由推导~$\rho_1$ 可得 $\Th{Q} \Proves
\ORefut(\num{m}, \gn{!R})$。由于我们假设 $\Th{T}$ 是一致的，故 $\Th{T} \Proves/ !R$。所以对所有 $k$，$k$ 都不是~$!R$ 的推导的Gödel数，因此由推导~$\pi_k$ 可得 $\Th{Q} \Proves \lnot
\OPrf(\num{k}, \gn{!R})$。于是我们有：
  
\begin{prooftree}\footnotesize
  \AxiomC{}
  \RightLabel{$\lambda_2$}
  \DeduceC{$\begin{gathered}[b] a = \num{0} \lor \dots \lor\\
      a = \num{m} \lor \num{m} < a\end{gathered}$}
  \AxiomC{$\dots$}
  \AxiomC{$\begin{gathered}\Discharge{\OPrf(a, \gn{!R})}{1}\\
    \Discharge{a = \num{k}}{2}\end{gathered}$}
  \RightLabel{$\pi_k'$}
  \DeduceC{$\lfalse$}
  \RightLabel{$\lfalse_I$}
  \UnaryInfC{$\num{m}<a$}
  \AxiomC{$\dots$}
  \AxiomC{$\Discharge{\num{m} < a}{2}$}
  \DischargeRule{$\Elim\lor^*$}{2}
  \insertBetweenHyps{\hspace{-.1em}}
  \doubleLine
  \QuinaryInfC{$\num{m} < a$}
  \AxiomC{}
  \RightLabel{$\rho_1$}
  \DeduceC{$\ORefut(\num{m}, \gn{!R})$}
  \RightLabel{\Intro\land}
    \insertBetweenHyps{\hspace{-.1em}}
  \BinaryInfC{$\num{m}
    < a \land \ORefut(\num{m}, \gn{!R})$}
  \RightLabel{\Intro\exists}
  \UnaryInfC{$\exists z(z
    < a \land \ORefut(z, \gn{!R}))$}
  \DischargeRule{\Intro\lif}{1}
  \UnaryInfC{$\OPrf(a, \gn{!R}) \lif \exists z(z
    < a \land \ORefut(z, \gn{!R}))$}
  \RightLabel{\Intro\forall}
  \UnaryInfC{$\forall x(\OPrf(x, \gn{!R}) \lif \exists z(z
    < x \land \ORefut(z, \gn{!R})))$}
  \DeduceC{$\lnot\exists x(\OPrf(x, \gn{!R}) \land \forall z(z
    < x \lif \lnot \ORefut(z, \gn{!R})))$}
\end{prooftree}
其中 $\pi_k'$ 是如下推导
\begin{prooftree}
  \AxiomC{}
  \RightLabel{$\pi_k$}
  \DeduceC{$\lnot \OPrf(\num{k}, \gn{!R})$}
  \AxiomC{$a = \num{k}$}
  \AxiomC{$\OPrf(a, \gn{!R})$}
  \RightLabel{\Elim=}
  \BinaryInfC{$\OPrf(\num{k}, \gn{!R})$}
  \RightLabel{\Elim\lnot}
  \BinaryInfC{$\lfalse$}
\end{prooftree}
而 $\lambda_2$ 是
\begin{prooftree}\footnotesize
  \AxiomC{}
  \RightLabel{$\lambda_3$}
  \DeduceC{$\begin{gathered}[b](a < \num{m} \lor {}\\a = \num{m}) \lor {}\\ \num{m} < a\end{gathered}$}

  \AxiomC{$\Discharge{a < \num{m}}{3}$}
  \RightLabel{$\lambda_1$}
  \DeduceC{$\begin{gathered}[b]a = \num{0} \lor \dots \lor {}\\
  a = \num{m-1}\end{gathered}$}
  %\RightLabel{\Intro\lor}
  \doubleLine
  \UnaryInfC{$\begin{gathered}[b]
      a = \num{0} \lor \dots \lor \\
      a = \num{m} \lor \num{m} < a
    \end{gathered}$}

  \AxiomC{$\Discharge{a = \num{m}}{3}$}
  %\RightLabel{\Intro\lor}
  \doubleLine
  \UnaryInfC{$\begin{gathered}[b]
      a = \num{0} \lor \dots \lor \\
      a = \num{m} \lor \num{m} < a
    \end{gathered}$}
  
  \AxiomC{$\Discharge{\num{m} < a}{3}$}
  \doubleLine
  \RightLabel{$\Intro\lor^*$}
  \UnaryInfC{$\begin{gathered}[b]
      a = \num{0} \lor \dots \lor \\
      a = \num{m} \lor \num{m} < a
    \end{gathered}$}

  %\insertBetweenHyps{\hskip 5pt}
  \doubleLine
  \DischargeRule{$\Elim\lor^2$}{3}
  \QuaternaryInfC{$a = \num{0} \lor \dots \lor a = \num{m} \lor \num{m} < a$}
\end{prooftree}
（推导 $\lambda_3$ 的存在性由 \olref[inc][req][min]{lem:trichotomy} 保证。我们用 $\Intro\lor^*$ 缩写重复使用 $\Intro\lor$，并将从 $(a < \num{m} \lor a = \num{m}) \lor \num{m} < a$ 推出 $a = \num{0} \lor \dots \lor a = \num{m} \lor \num{m} < a$ 时双重使用 $\Elim\lor$ 记为 $\Elim\lor^2$。）
\end{proof}

\OLEndChapterHook


\let\intro\comment
\let\endintro\endcomment

\chapter{一阶逻辑}\label{fol:chap}

\olimport*[first-order-logic/syntax-and-semantics]{first-order-languages}

\olimport*[first-order-logic/syntax-and-semantics]{terms-formulas}

\olimport*[first-order-logic/syntax-and-semantics]{free-vars-sentences}

\olimport*[first-order-logic/syntax-and-semantics]{substitution}

\olimport*[first-order-logic/syntax-and-semantics]{structures}

\olimport*[first-order-logic/syntax-and-semantics]{satisfaction}

\olimport*[first-order-logic/syntax-and-semantics]{assignments}


\olimport*[first-order-logic/syntax-and-semantics]{extensionality}

\olimport*[first-order-logic/syntax-and-semantics]{semantic-notions}

\section{理论}

\begin{defn}
一个语句集~$\Gamma$ 是\emph{封闭的}，当且仅当：只要 $\Gamma \Entails !A$，就有 $!A \in \Gamma$。语句集~$\Gamma$ 的\emph{闭包}是 $\Setabs{!A}{\Gamma \Entails !A}$。

我们说~$\Gamma$ \emph{由}语句集~$\Delta$ \emph{公理化}，是指 $\Gamma$ 是~$\Delta$ 的闭包。
\end{defn}

\begin{ex}
语言~$\Lang L_<$ 中的严格线性序理论由集合
\begin{align*}
& \lforall[x][\lnot x < x], \\
& \lforall[x][\lforall[y][((x < y \lor y <
    x) \lor x = y)]], \\
& \lforall[x][\lforall[y][\lforall[z][((x < y
      \land y < z) \lif x < z)]]]
\end{align*}
公理化。它完全刻画了预期的结构：每个严格线性序都是此公理系统的模型，反之，若 $R$ 是集合 $X$ 上的线性序，则结构 $\Struct M$（其中 $\Domain M = X$ 且 $\Assign{<}{M} = R$）就是这个理论的模型。
\end{ex}

\begin{ex}
语言 $\Obj 1$（常项符号）、$\cdot$（二元函数符号）中的群理论由
\begin{align*}
& \lforall[x][(x \cdot \Obj 1) = x]\\
& \lforall[x][\lforall[y][\lforall[z][\eq[(x \cdot (y \cdot z))][((x
          \cdot y) \cdot z)]]]]\\
& \lforall[x][\lexists[y][(x \cdot y) = \Obj 1]]
\end{align*}
公理化。
\end{ex}

\begin{ex}
Peano（皮亚诺）算术理论由算术语言~$\Lang L_A$ 中的下列语句公理化：
\begin{align*}
& \lnot\lexists[x][x' = \Obj 0]\\
& \lforall[x][\lforall[y][(x' = y' \lif x = y)]]\\
& \lforall[x][\lforall[y][(x < y \liff \lexists[z][\eq[(z' + x)][y]])]]\\
& \lforall[x][\eq[(x + \Obj 0)][x]]\\
& \lforall[x][\lforall[y][\eq[(x + y')][(x + y)']]]\\
& \lforall[x][\eq[(x \times \Obj 0)][\Obj 0]]\\
& \lforall[x][\lforall[y][\eq[(x \times y')][((x \times y) + x)]]]\\
\intertext{再加上所有如下形式的句子}
& (!A(\Obj 0) \land \lforall[x][(!A(x) \lif !A(x'))]) \lif \lforall[x][!A(x)]
\end{align*}
由于后一种形式有无穷多语句，这个公理系统是无限的。后一种形式称为\emph{归纳模式}。（实际上，归纳模式比我们这里描述的要稍微复杂一点。）

第三条公理是~$<$ 的一个\emph{显定义}。
\end{ex}

\OLEndChapterHook

\chapter{自然演绎}\label{nd:chap}

\olimport*[first-order-logic/proof-systems]{natural-deduction}[\nosection]

\olimport*[first-order-logic/natural-deduction]{rules-and-proofs}

\olimport*[first-order-logic/natural-deduction]{propositional-rules}

\olimport*[first-order-logic/natural-deduction]{quantifier-rules}

\olimport*[first-order-logic/natural-deduction]{derivations}

\olimport*[first-order-logic/natural-deduction]{proving-things}

\olimport*[first-order-logic/natural-deduction]{proving-things-quant}

\olimport*[first-order-logic/natural-deduction]{identity}

\olimport*[first-order-logic/natural-deduction]{proof-theoretic-notions}

\olimport*[first-order-logic/natural-deduction]{soundness}

\OLEndChapterHook

%\olimport[additions]{computing-recursive}

\stopproblems
\def\ifproblems#1{}

\def\figurename{图}

\chapter{传记}\label{bios:chap}

\olimport*[history/biographies]{alonzo-church}

\olimport*[history/biographies]{kurt-goedel}

\olimport*[history/biographies]{rozsa-peter}

\olimport*[history/biographies]{julia-robinson}

\olimport*[history/biographies]{alfred-tarski}

\backmatter

\clearpage
\photocredits

\bibliographystyle{\olpath/bib/natbib-oup}
\bibliography{\olpath/bib/open-logic.bib}

\olimport*{\olpath/content/open-logic-about}